Let \(\mathcal{F}\subseteq 2^{[n]}\) be an antichain.
Fix \(r\in\mathbb{N}\) and assume that for every size \(k\) occurring in \(\mathcal{F}\), there are at least \(r\) sets in \(\mathcal{F}\) of size \(k\).
How large can \(|\{|A|:A\in\mathcal{F}\}|\) be?
For \(r=1\) the maximum is \(n-2\).
For \(r>1\) it is stated that \(n-2\) is impossible but \(n-3\) is achievable for large \(n\).
Question: determine/estimate how large \(n\) must be (as a function of \(r\)) to guarantee an example with \(n-3\) distinct sizes.
